% Seção I — Introdução com subseções obrigatórias
\section{Introdução}
\label{sec:introducao}

[Parágrafo introdutório contextualizando o tema e apresentando o problema de forma geral.
Deve situar o leitor no domínio do trabalho e motivar a leitura das seções seguintes.
Aproximadamente 3 parágrafos.]

\subsection{Justificativa}
\label{sec:justificativa}

[Por que este trabalho é relevante? Qual o impacto prático e acadêmico?
Mencionar demanda da área, lacuna identificada na literatura ou necessidade social.
Aproximadamente 2 parágrafos.]

\subsection{Objetivos}
\label{sec:objetivos}

\subsubsection{Objetivo Geral}
\label{sec:objetivo-geral}

[Uma frase clara e direta descrevendo o objetivo principal do trabalho, iniciando com verbo no infinitivo.]

\subsubsection{Objetivos Específicos}
\label{sec:objetivos-especificos}

\begin{itemize}
  \item [Objetivo específico 1 --- verbo no infinitivo];
  \item [Objetivo específico 2 --- verbo no infinitivo];
  \item [Objetivo específico 3 --- verbo no infinitivo];
  \item [Objetivo específico 4 --- verbo no infinitivo].
\end{itemize}

\subsection{Estrutura do Trabalho}
\label{sec:estrutura}

[Descrever brevemente o conteúdo de cada seção subsequente.
Exemplo: ``A Seção~\ref{sec:referencial} apresenta o referencial teórico que fundamenta este trabalho.
A Seção~\ref{sec:metodologia} descreve a metodologia adotada para o desenvolvimento da solução...'']
